\subsection{对偶化模的结构}

引理 2.—— 设 A 为 Noether 环，M 与 N 为有限生成的 A-模，\(u: \mathrm{M} \to \mathrm{N}\) 为一个同态。设 \(x\) 为 A 的根的一个元素，使同态 \(x_{\mathrm{N}}\) 是单射的。若同态 \(\overline{u}: \mathrm{M}/x\mathrm{M} \to \mathrm{N}/x\mathrm{N}\)（它由 \(u\) 诱导）是单射的（相应地，满射的、双射的），则 \(u\) 亦然。

关于 \(u\) 的满射性的断言由 Nakayama 引理（II, § 3, 第 2 节，命题 4 的推论 1）得出，无须对 \(x_{\mathrm{N}}\) 作任何假设。考虑行正合的交换图表

\[
\begin{array}{c c c c c c c} & \mathrm{M} & \xrightarrow {x _ {\mathrm{M}}} & \mathrm{M} & \longrightarrow & \mathrm{M} / x \mathrm{M} \\ & u \Bigg \downarrow & & u \Bigg \downarrow & & \overline {{u}} \Bigg \downarrow \\ 0 & \longrightarrow & \mathrm{N} & \xrightarrow {x _ {\mathrm{N}}} & \mathrm{N} & \longrightarrow & \mathrm{N} / x \mathrm{N} \end{array} ;
\]

利用蛇形引理（I, § 1, 第 4 节，命题 2），可导出一个正合列 \(\operatorname{Ker} u \xrightarrow{x} \operatorname{Ker} u \longrightarrow \operatorname{Ker} \overline{u}\)。若 \(\overline{u}\) 是单射的，则以 \(x\) 为比的乘法映射在 \(\operatorname{Ker} u\) 上是满射的，由 Nakayama 引理这蕴含 \(\operatorname{Ker} u = 0\)。

命题 7.—— 设 A 为 Noether 环，Ω 为对偶化 A-模。

\begin{enumerate}
\def\labelenumi{\alph{enumi})}
\item
  有 \(\operatorname{Ext}_{\mathrm{A}}^{i}(\Omega, \Omega) = 0\)（对 \(i > 0\)）。
\item
  \(L'\) 典范同态 \(\gamma : \mathrm{A} \to \operatorname{End}_{\mathrm{A}}(\Omega)\) 是双射的。
\item
  每个对偶化 A-模都具有形状 \(\Omega\otimes_{\mathrm{A}}\mathrm{L}\)，其中 \(\mathrm{L}\) 是秩为 1 的投射 A-模。
\end{enumerate}

\begin{enumerate}
\def\labelenumi{\Alph{enumi})}
\tightlist
\item
  先处理环 A 是局部环的情形。在此情形，条件 c) 只是说两个对偶化模是同构的。
\end{enumerate}

设 \(\Omega'\) 为对偶化 A-模。有 \(\operatorname{prof}_{\mathrm{A}}(\Omega') = \dim_{\mathrm{A}}(\Omega') = \dim(\mathrm{A})\)（第 1 节，命题 1），故有 \(\operatorname{Ext}_{\Lambda}^{i}(\Omega', \Omega) = 0\)（对 \(i \neq 0\)）（第 1 节，命题 3, c)），由此得 a)。

我们对整数 dim(A)（它等于 prof(A)）作归纳，来证明 b) 与 c)。若它为零，则环 A 是 Artin 环，\(\Omega'\) 与 \(\Omega\) 都是 Matlis A-模（\(n^{\circ}\) 1, 例 1）；因而它们是同构的（\(\S\) 8, \(n^{\circ}\)），且典范映射 A → \(\mathrm{End}_{\mathrm{A}}(\Omega)\) 是双射的（\(\S\) 8, \(n^{\circ}\) 2, 命题 3, c)）。设 dim(A) \textgreater{} 0，并设 \(x\) 为 \(\mathfrak{m}_{\mathrm{A}}\) 的一个可消去元素。同态 \(x_{\Omega}\) 是单射的（\(n^{\circ}\) 2, 命题 4），并且有一个正合列

\[
0 \rightarrow \Omega \xrightarrow {x _ {\Omega}} \Omega \longrightarrow \Omega / x \Omega \rightarrow 0.
\]

由于 \(\mathrm{Ext}_{\mathrm{A}}^{1}(\Omega',\Omega)\) 为零，且 \(\mathrm{Hom}_{\mathrm{A}}(\Omega',\Omega /x\Omega)\) 等同于 \(\mathrm{Hom}_{\mathrm{A} / x\mathrm{A}}(\Omega'/x\Omega',\Omega/x\Omega)\)，可导出一个正合列

\[
0 \rightarrow \operatorname{Hom} _ {\mathrm{A}} \left(\Omega^ {\prime}, \Omega\right) \xrightarrow {x} \operatorname{Hom} _ {\mathrm{A}} \left(\Omega^ {\prime}, \Omega\right) \xrightarrow {p} \operatorname{Hom} _ {\mathrm{A} / x \mathrm{A}} \left(\Omega^ {\prime} / x \Omega^ {\prime}, \Omega / x \Omega\right)\rightarrow 0,\tag{1}
\]

其中 \(p\) 是典范映射。由命题 4，A/xA-模 \(\Omega/x\Omega\) 与 \(\Omega'/x\Omega'\) 都是对偶化的，因而由归纳假设它们同构。设 \(\overline{u}\) 是从 \(\Omega'/x\Omega'\) 到 \(\Omega/x\Omega\) 的一个同构。考虑到正合列 (1)，存在一个 A-同态 \(u: \Omega' \to \Omega\)，使 \(p(u) = \overline{u}\)；由引理 2，\(u\) 是双射的，这就证明了 c)。由归纳假设，典范同态 A/xA \(\longrightarrow\) End \(_{A/xA}(\Omega/x\Omega)\) 是双射的。鉴于正合列 (1)，此同态等同于同态 \(\overline{\gamma}: A/xA \longrightarrow \text{End}_A(\Omega)/x \text{End}_A(\Omega)\)（它由 \(\gamma\) 诱导）；于是由引理 2 推出 \(\gamma\) 是双射的，由此得 b)。

\begin{enumerate}
\def\labelenumi{\Alph{enumi})}
\setcounter{enumi}{1}
\tightlist
\item
  转到一般情形。对 A 的每个极大理想 m 及每个整数 i \textgreater{} 0，由前所述有 \(\operatorname{Ext}_{A_{m}}^{i}(\Omega_{\mathfrak{m}}, \Omega_{\mathfrak{m}}) = 0\)，因此 \(\operatorname{Ext}_{A}^{i}(\Omega, \Omega)_{\mathfrak{m}} = 0\)（§ 3, 第 2 节，命题 2），这蕴含 \(\operatorname{Ext}_{A}^{i}(\Omega, \Omega) = 0\)（II, § 3, 第 3 节，定理 1 的推论 2）。同样，同态 \(\gamma_{m}: A_{m} \to \operatorname{End}_{A}(\Omega)_{m}\) 对 A 的每个极大理想 m 都是双射的，故 \(\gamma\) 是双射的（前引文，定理 1）。
\end{enumerate}

最后证明 c)。设 \(\Omega'\) 为对偶化 A-模。用 L 表示 A-模 \(\mathrm{Hom}_{\mathrm{A}}(\Omega',\Omega)\)，并用 \(v:\Omega' \otimes_{\mathrm{A}} \mathrm{L} \longrightarrow \Omega\) 表示满足 \(v(x \otimes f) = f(x)\)（对 \(x \in \Omega'\)，\(f \in \mathrm{L}\)）的同态。设 \(\mathfrak{m}\) 为 A 的一个极大理想。\(\mathrm{A}_{\mathfrak{m}}\)-模 \(\mathrm{L}_{\mathfrak{m}}\) 等同于 \(\mathrm{Hom}_{\mathrm{A}_{\mathfrak{m}}}(\Omega_{\mathfrak{m}}',\Omega_{\mathfrak{m}})\)；由已处理的情形，它是秩为 1 的自由模，且每个同构 \(h:\Omega_{\mathfrak{m}}' \to \Omega_{\mathfrak{m}}\) 都是它的一个生成元。当利用生成元 \(\mathrm{L}_{\mathfrak{m}}\) 把 \(\mathrm{A}_{\mathfrak{m}}\) 与 \(h\) 等同时，同态 \(v_{\mathfrak{m}}:\Omega_{\mathfrak{m}}' \otimes_{\mathrm{A}_{\mathfrak{m}}}\mathrm{L}_{\mathfrak{m}} \longrightarrow \Omega_{\mathfrak{m}}\) 等同于 \(h\)，因此它是双射的。由于这对 A 的每个极大理想 \(\mathfrak{m}\) 都成立，A-模 L 是秩为 1 的投射模（II, § 5, 第 3 节，定理 2），且同态 \(v\) 是双射的（II, § 3, 第 3 节，定理 1）。

推论 1.—— 为使 A 是 Gorenstein 环，必须且只须 A-模 Ω 是秩为 1 的投射模。

这由第 1 节的例 2 及命题 7 c) 得出。

推论 2.—— 设环 A 是可表现的。A 的使 A\(_{p}\) 为 Gorenstein 环的素理想 p 的集合在 Spec(A) 中是开的。

设 p 为 A 的一个使 \(A_{p}\) 是 Gorenstein 环的素理想。于是它是 Macaulay 环；必要时把 A 换成 \(A_{f}\)，其中 f 是 A − p 的一个适当元素，就可化归为 A 是 Macaulay 环的情形（§ 4, 第 4 节，命题 7, c)）。设 \(\Omega\) 为对偶化 A-模（第 3 节，命题 6 的推论 3）。于是 \(\Omega_{p}\) 是 Gorenstein 环 \(A_{p}\) 上的对偶化模（第 1 节，命题 2），从而是秩为 1 的自由模（推论 1）。因此存在 A − p 的一个元素 g，使 \(A_{g}\)-模 \(\Omega_{g}\) 是秩为 1 的自由模（II, § 5, 第 1 节，命题 2 的推论）。于是 \(A_{g}\) 是 Gorenstein 环（推论 1），且对 A 的每个不含 g 的素理想 q，\(A_{q}\) 亦然（§ 3, 第 7 节，例 1），这就证明了推论。

